\documentclass[10pt,a4paper]{amsart}
\usepackage[margin=19mm]{geometry}
\usepackage{amsmath,amssymb,mathtools}
\usepackage[scr=rsfs,cal=euler]{mathalfa}
\usepackage{xfrac}
\usepackage[unicode,hidelinks,bookmarks=false]{hyperref}
\newcommand{\ie}{i.e.\ }
\newcommand{\cf}{cf.\ }
\newcommand{\resp}{respectively\ }
\newcommand{\Subsection}{\subsection}
\providecommand{\nobreakdash}{\nobreak\hbox{-}}
\setlength{\emergencystretch}{2em}
\multlinegap0pt
\usepackage{amssymb}
\usepackage{mathtools}
\usepackage{xfrac}
\usepackage{bm}
\usepackage{algorithm}\usepackage{algpseudocode}
\usepackage{xcolor}
\newtheorem{lemma}{Lemma}
\usepackage{cleveref}
\crefname{lemma}{Lemma}{Lemmas}
\newcommand{\KL}[2]{D_{\mathrm{KL}}\left(#1 \,\|\, #2\right)} 
\newcommand{\R}{\mathbb{R}}
\newcommand{\Tr}{\operatorname{Tr}}
\newcommand{\bAblk}[1]{\mathbf{A}_{[#1]}}
\newcommand{\bG}{\mathbf{G}}
\newcommand{\bX}{\mathbf{X}}
\newcommand{\bY}{\mathbf{Y}}
\newcommand{\bw}{\mathbf{w}}
\newcommand{\bwbar}{\overline{\mathbf{w}}}
\newcommand{\bx}{\mathbf{x}}
\newcommand{\by}{\mathbf{y}}
\newcommand{\calT}{\mathcal{T}}
\newcommand{\kldiv}[2]{\mathcal{D}_{\mathrm{KL}}\left(#1 \,\|\, #2\right)}
\newcommand\numberthis{\addtocounter{equation}{1}\tag{\theequation}}
\newcommand{\omrp}{\gamma} 
\newcommand{\one}{\mathbf{1}}
\newcommand{\wbar}{\overline{w}}
\title{Computing Lewis Weights to High Precision by Fixed-Point Iteration}
\author{Swati Padmanabhan}
\date{Source: arXiv:2609.04338v1 (CC BY 4.0)}
\begin{document}
\maketitle

\noindent\emph{Source and licence.} Computing Lewis Weights to High Precision by Fixed-Point Iteration. Version arXiv:2609.04338v1 (CC BY 4.0), distributed by arXiv under the Creative Commons Attribution 4.0 International licence, http://creativecommons.org/licenses/by/4.0/, as recorded in the arXiv licence metadata for this exact version and shown on the abstract page https://arxiv.org/abs/2609.04338v1. Full licence notice: \texttt{licenses/arxiv-2609.04338-CC-BY-4.0.txt}.\par\smallskip

\noindent\textit{Editorial scope.} Counted unit: the article's lemma lem:kl-contraction (source0.tex lines 545--550). The article's complete original proof of it is delivered with it (source0.tex lines 552--559, one \texttt{proof} environment of 8 lines). No mathematics is written, repaired or re-derived: the statement and the proof are the article's own lines, in the article's own notation, and no retained line is substituted or re-flowed.  Retained uncounted as the source-local context the proof needs: lines 514--521.  Nothing was omitted from any retained span, so the package carries no omission record.  No retained line contains a citation, so the wrapper carries no bibliography and no bibliography entry.  No retained line contains a figure environment, an image-inclusion command or drawing code, so no image is delivered and the package declares no asset dependency.  Editorial formatting only, in addition: an AMS \texttt{amsart} preamble that restates the article's own declarations and macros used by the retained lines, namely the environments \texttt{lemma} on separate counters, together with the standard packages those lines need.  The retained environments print in this document as Lemma 1 in order of appearance, whereas the article prints the same items with the numbering of its own full document.  The complete native source of the article as distributed in the arXiv e-print for this exact version is bundled unchanged as \texttt{sources/209360/source0.tex}.
More generally, for positive vectors $\bx,\by\in\R_{>0}^d$ of the same dimension satisfying $\one^\top\bx=\one^\top\by=n$, we define the scaled Kullback-Leibler (KL) divergence
\[
\kldiv{\bx}{\by}
:=
\sum_{i=1}^d x_i\log\left(\frac{x_i}{y_i}\right).
\numberthis\label{eq:notation-KLdiv}
\]
For probability distributions $\mu$ and $\nu$, we use $\KL{\mu}{\nu}$ for the usual KL divergence.

\begin{lemma}
\label{lem:kl-contraction}
Let  $\bw\in\Omega$. Then  $\kldiv{\bwbar}{\calT(\bw)}
    \leq
    \omrp\kldiv{\bwbar}{\bw}.$
\end{lemma}

\begin{proof}
By  the definition of block leverage scores and $\bwbar=\calT(\bwbar)$, we have \[ \calT_i(\bw)=w_i^\omrp \Tr\left(\bAblk{i}\bG(\bw)^{-1}\bAblk{i}^\top\right), \qquad \wbar_i = \wbar_i^{\omrp} \Tr\left(\bAblk{i}\bG(\bwbar)^{-1}\bAblk{i}^\top\right).\] Therefore, by expanding out the definition in \Cref{eq:notation-KLdiv}, we have \[\kldiv{\bwbar}{\calT(\bw)} = \omrp \kldiv{\bwbar}{\bw} + \sum_{i = 1}^k \wbar_i \log\left(\frac{\Tr\left(\bAblk{i}\bG(\bwbar)^{-1}\bAblk{i}^\top\right)}{\Tr\left(\bAblk{i}\bG(\bw)^{-1}\bAblk{i}^\top\right)}\right). \numberthis\label{eq:KLcon-2}\] We  complete the  proof  by showing next that the second term of \Cref{eq:KLcon-2} is nonpositive. To this end, we first note that by Jensen's inequality applied to $x\mapsto\log(x)$ and the fact that $\sum_{i=1}^k \wbar_i = n$, we have  \[ \sum_{i = 1}^k \frac{\wbar_i}{n} \log\left(\frac{\Tr\left(\bAblk{i}\bG(\bwbar)^{-1}\bAblk{i}^\top\right)}{\Tr\left(\bAblk{i}\bG(\bw)^{-1}\bAblk{i}^\top\right)}\right) \leq \log\left( \sum_{i=1}^k \frac{\wbar_i}{n} \frac{\Tr\left(\bAblk{i}\bG(\bwbar)^{-1}\bAblk{i}^\top\right)}{\Tr\left(\bAblk{i}\bG(\bw)^{-1}\bAblk{i}^\top\right)} \right),\] and it  suffices to show that \[\sum_{i=1}^k {\wbar_i} \frac{\Tr\left(\bAblk{i}\bG(\bwbar)^{-1}\bAblk{i}^\top\right)}{\Tr\left(\bAblk{i}\bG(\bw)^{-1}\bAblk{i}^\top\right)} \stackrel{?}{\leq} n.\numberthis\label{eq:KLcon-3}\]To this end, we use $\bwbar=\calT(\bwbar)$ and apply matrix Cauchy-Schwarz inequality $\left|\Tr\left(\bX^\top\bY\right)\right|\leq \|\bX\|_{\mathrm{F}}\|\bY\|_{\mathrm{F}}$ with $\bX:=\bG(\bw)^{1/2}\bG(\bwbar)^{-1}\bAblk{i}^\top$ and $\bY:= \bG(\bw)^{-1/2}\bAblk{i}^\top$ and first obtain \[\wbar_i \frac{\Tr\left(\bAblk{i} \bG(\bwbar)^{-1}\bAblk{i}^\top\right)}{\Tr\left(\bAblk{i} \bG(\bw)^{-1}\bAblk{i}^\top\right)} = \wbar_i^{\omrp} \frac{\left[\Tr \left(\bAblk{i}\bG(\bwbar)^{-1}\bAblk{i}^\top\right)\right]^2}{\Tr\left(\bAblk{i}\bG(\bw)^{-1}\bAblk{i}^\top\right)} \leq \wbar_i^{\omrp} \Tr\left(\bAblk{i} \bG(\bwbar)^{-1}\bG(\bw)  \bG(\bwbar)^{-1} \bAblk{i}^\top\right).\] Summing the above inequality over $i = 1, 2, \dotsc, k$ and applying the definition of $\bG(\bwbar)$ and $\bG(\bw)$, linearity of trace, and the invariance of trace to cyclic permutation, we have
\begin{align*}
\sum_{i=1}^k \wbar_i\frac{\Tr\left(\bAblk{i}\bG(\bwbar)^{-1}\bAblk{i}^\top\right)}{\Tr\left(\bAblk{i}\bG(\bw)^{-1}\bAblk{i}^\top\right)} &\leq \sum_{i=1}^k \wbar_i^\omrp \Tr\left(\bAblk{i}\bG(\bwbar)^{-1}\bG(\bw) \bG(\bwbar)^{-1} \bAblk{i}^\top\right)= \Tr\left(\bG(\bw)\bG(\bwbar)^{-1}\right).\numberthis\label{eq:KLcon-4}
\end{align*}
Next, again by $\bwbar=\calT(\bwbar)$, H\"older's inequality with exponents $\frac{1}{\omrp}$ and $\frac{1}{1-\omrp}$ (valid since we restrict ourselves to $p>2$), and the fact that $\bw\in\Omega$, we have \[ \Tr\left(\bG(\bw)\bG(\bwbar)^{-1}\right)
= \sum_{i=1}^k w_i^\omrp \Tr\left( \bAblk{i}\bG(\bwbar)^{-1} \bAblk{i}^\top\right) = \sum_{i=1}^k w_i^\omrp \wbar_i^{1-\omrp}\leq \left(\sum_{i=1}^k w_i\right)^\omrp \left(\sum_{i=1}^k \wbar_i\right)^{1-\omrp} = n. \numberthis\label{eq:KLcon-5}\]  Combining \Cref{eq:KLcon-4} and \Cref{eq:KLcon-5} completes the proof of \Cref{eq:KLcon-3}, which finishes the proof of the claim. 
\end{proof}

\end{document}
